% ============================================================
%  LaTeX source: Section 2 & Section 3
%  Paper: "Error estimates for the finite element method of the
%          Navier-Stokes-Poisson-Nernst-Planck equations"
%  Authors: Minghao Li, Zhenzhen Li
%  Journal: Applied Numerical Mathematics 197 (2024) 186-209
% ============================================================
%
% Compile with: pdflatex -> bibtex -> pdflatex -> pdflatex
%
% Required packages (add to preamble):
%   \usepackage{amsmath, amssymb, amsthm}
%   \usepackage{mathtools}
%
% Theorem environments assumed defined as:
%   \newtheorem{lemma}{Lemma}
%   \newtheorem{theorem}{Theorem}
%   \newtheorem{remark}{Remark}
% ============================================================

% ------------------------------------------------------------------
\section{Preliminary}
\label{sec:preliminary}
% ------------------------------------------------------------------

We introduce the following function spaces
\begin{equation}
  X = H^{1}(\Omega),\quad
  X^{0} = \Bigl\{\,\eta\in H^{1}(\Omega) : \int_{\Omega}\eta\,\mathrm{d}x = 0\,\Bigr\},
  \label{eq:X-space}
\end{equation}
\begin{equation}
  V = \bigl(H^{1}_{0}(\Omega)\bigr)^{d},\quad
  Q = L^{2}_{0}(\Omega) = \Bigl\{\,q\in L^{2}(\Omega) : \int_{\Omega}q\,\mathrm{d}x = 0\,\Bigr\}.
  \label{eq:VQ-space}
\end{equation}

The trilinear forms $b_{1}(\cdot\,;\cdot,\cdot)$ on $V\times V\times V$ and
$b_{2}(\cdot\,;\cdot,\cdot)$ on $V\times X\times X$ are defined by
\begin{equation}
  b_{1}(u;\,v,\,w)
  = \tfrac{1}{2}\bigl((u\cdot\nabla)v,\,w\bigr)
  - \tfrac{1}{2}\bigl((u\cdot\nabla)w,\,v\bigr),
  \quad \forall\,u,v,w\in V,
  \label{eq:b1}
\end{equation}
\begin{equation}
  b_{2}(u;\,c,\,\eta)
  = \tfrac{1}{2}(u\cdot\nabla c,\,\eta)
  - \tfrac{1}{2}(u\cdot\nabla\eta,\,c),
  \quad \forall\,u\in V,\;c,\eta\in X.
  \label{eq:b2}
\end{equation}

We recall the following Poincar\'{e}--Friedrichs inequalities
\begin{equation}
  \Vert v\Vert_{1} \le C\,\Vert\nabla v\Vert_{0},
  \quad \forall\,v\in V,
  \label{eq:PF1}
\end{equation}
\begin{equation}
  \Vert\varphi\Vert_{1}
  \le C\!\left(\Vert\nabla\varphi\Vert_{0}
    + \Bigl|\int_{\Omega}\varphi\,\mathrm{d}x\Bigr|\right),
  \quad \forall\,\varphi\in X,
  \label{eq:PF2}
\end{equation}
and the embedding inequalities
\begin{equation}
  \Vert\varphi\Vert_{0,6} \le C\,\Vert\varphi\Vert_{1},
  \quad \forall\,\varphi\in X.
  \label{eq:embed}
\end{equation}

Next, integrating the third equation of \eqref{eq:NSPNP} and using the
Green formula, we can derive
\begin{equation}
  \frac{\mathrm{d}}{\mathrm{d}t}
  \!\left(\int_{\Omega}c_{1}\,\mathrm{d}x\right)
  - \int_{\partial\Omega}\frac{\partial c_{1}}{\partial n}\,\mathrm{d}s
  - \int_{\partial\Omega}c_{1}\frac{\partial\psi}{\partial n}\,\mathrm{d}s
  + \int_{\partial\Omega}u\cdot n\,c_{1}\,\mathrm{d}s = 0.
  \label{eq:c1-int1}
\end{equation}
Then in view of the boundary conditions, we get
\begin{equation}
  \frac{\mathrm{d}}{\mathrm{d}t}
  \!\left(\int_{\Omega}c_{1}\,\mathrm{d}x\right) = 0,
  \label{eq:c1-conserve}
\end{equation}
which implies
\begin{equation}
  \int_{\Omega}c_{1}\,\mathrm{d}x = \int_{\Omega}c_{10}\,\mathrm{d}x.
  \label{eq:c1-mass}
\end{equation}
Manipulating the fourth and fifth equations of \eqref{eq:NSPNP} in the
same way, we obtain
\begin{equation}
  \int_{\Omega}c_{2}\,\mathrm{d}x = \int_{\Omega}c_{20}\,\mathrm{d}x,
  \qquad
  \int_{\Omega}c_{1}\,\mathrm{d}x = \int_{\Omega}c_{2}\,\mathrm{d}x.
  \label{eq:c12-mass}
\end{equation}
Set
\begin{equation}
  m_{0}
  = \frac{1}{|\Omega|}\int_{\Omega}c_{10}\,\mathrm{d}x
  = \frac{1}{|\Omega|}\int_{\Omega}c_{20}\,\mathrm{d}x.
  \label{eq:m0}
\end{equation}

Take $\tilde{c}_{1}=c_{1}-m_{0}$, $\tilde{c}_{2}=c_{2}-m_{0}$; then we
obtain the variational form of equations \eqref{eq:NSPNP}: find
$(u,p,\tilde{c}_{1},\tilde{c}_{2},\psi)\in V\times Q\times X^{0}\times X^{0}\times X^{0}$
such that
\begin{equation}
  \left\{
  \begin{aligned}
    &(u_{t},v) + (\nabla u,\nabla v) + b_{1}(u;u,v)
      - (\operatorname{div}v,\,p)
      = -\bigl((\tilde{c}_{1}-\tilde{c}_{2})\nabla\psi,\,v\bigr),
      &&\forall\,v\in V,\\[4pt]
    &(\operatorname{div}u,\,q) = 0,
      &&\forall\,q\in Q,\\[4pt]
    &(\tilde{c}_{1t},\,\eta_{1})
      + (\nabla\tilde{c}_{1},\nabla\eta_{1})
      + \bigl((\tilde{c}_{1}+m_{0})\nabla\psi,\,\nabla\eta_{1}\bigr)
      + b_{2}(u;\tilde{c}_{1},\eta_{1}) = 0,
      &&\forall\,\eta_{1}\in X,\\[4pt]
    &(\tilde{c}_{2t},\,\eta_{2})
      + (\nabla\tilde{c}_{2},\nabla\eta_{2})
      - \bigl((\tilde{c}_{2}+m_{0})\nabla\psi,\,\nabla\eta_{2}\bigr)
      + b_{2}(u;\tilde{c}_{2},\eta_{2}) = 0,
      &&\forall\,\eta_{2}\in X,\\[4pt]
    &(\nabla\psi,\nabla\varphi) = (\tilde{c}_{1}-\tilde{c}_{2},\,\varphi),
      &&\forall\,\varphi\in X.
  \end{aligned}
  \right.
  \label{eq:variational}
\end{equation}

Let $\mathcal{T}_{h}$ be a regular and quasi-uniform triangular or
tetrahedral partition of $\Omega$. We consider low-order conforming finite
element spaces
$V_{h}\times Q_{h}\times X^{0}_{h}\times X^{0}_{h}\times X^{0}_{h}
\subset V\times Q\times X^{0}\times X^{0}\times X^{0}$,
which satisfy the following approximation properties: for
$u,\tilde{c}_{1},\tilde{c}_{2},\psi\in H^{2}(\Omega)$,
$p\in H^{1}(\Omega)$, there exist
$v_{h},q_{h},\chi_{1h},\chi_{2h},\varphi_{h}
\in V_{h}\times Q_{h}\times X^{0}_{h}\times X^{0}_{h}\times X^{0}_{h}$
such that
\begin{equation}
  \Vert\tilde{c}_{i}-\chi_{ih}\Vert_{0}
  + h\,\Vert\nabla(\tilde{c}_{i}-\chi_{ih})\Vert_{0}
  \le Ch^{2}\Vert\tilde{c}_{i}\Vert_{2},
  \quad i=1,2,
  \label{eq:approx-c}
\end{equation}
\begin{equation}
  \Vert\psi-\varphi_{h}\Vert_{0}
  + h\,\Vert\nabla(\psi-\varphi_{h})\Vert_{0}
  \le Ch^{2}\Vert\psi\Vert_{2},
  \label{eq:approx-psi}
\end{equation}
\begin{equation}
  \Vert u-v_{h}\Vert_{0}
  + h\bigl(\Vert\nabla(u-v_{h})\Vert_{0}
         + \Vert p-q_{h}\Vert_{0}\bigr)
  \le Ch^{2}\bigl(\Vert u\Vert_{2}+\Vert p\Vert_{1}\bigr),
  \label{eq:approx-u}
\end{equation}
and the inf-sup condition \cite{John2016}
\begin{equation}
  \sup_{v_{h}\in V_{h}}
  \frac{(\operatorname{div}v_{h},\,q_{h})}{\Vert\nabla v_{h}\Vert_{0}}
  \ge C\,\Vert q_{h}\Vert_{0},
  \quad \forall\,q_{h}\in Q_{h}.
  \label{eq:infsup}
\end{equation}

We define the Ritz projection
$R_{h}:X^{0}\to X^{0}_{h}$ and the Stokes projection
$S_{h}:V\times Q\to V_{h}\times Q_{h}$ by
\begin{equation}
  \bigl(\nabla(R_{h}\eta-\eta),\,\nabla\chi_{h}\bigr) = 0,
  \quad \forall\,\chi_{h}\in X^{0}_{h},
  \label{eq:Ritz}
\end{equation}
and
\begin{equation}
  \left\{
  \begin{aligned}
    &(\nabla S_{h}u,\nabla v_{h}) - (\operatorname{div}v_{h},S_{h}p)
      = (\nabla u,\nabla v_{h}) - (\operatorname{div}v_{h},p),
      &&\forall\,v_{h}\in V_{h},\\
    &(\operatorname{div}S_{h}u,\,q_{h}) = 0,
      &&\forall\,q_{h}\in Q_{h}.
  \end{aligned}
  \right.
  \label{eq:Stokes-proj}
\end{equation}

Then, from the standard finite element theory \cite{John2016,Li2021}, we have
\begin{equation}
  \Vert\eta - R_{h}\eta\Vert_{0}
  + h\,\Vert\nabla(\eta-R_{h}\eta)\Vert_{0}
  \le Ch^{2}\Vert\eta\Vert_{2},
  \label{eq:Ritz-est}
\end{equation}
\begin{equation}
  \Vert u-S_{h}u\Vert_{0}
  + h\bigl(\Vert\nabla(u-S_{h}u)\Vert_{0}
    + \Vert p-S_{h}p\Vert_{0}\bigr)
  \le Ch^{2}\bigl(\Vert u\Vert_{2}+\Vert p\Vert_{1}\bigr),
  \label{eq:Stokes-est}
\end{equation}
\begin{equation}
  \Vert u_{t} - S_{h}u_{t}\Vert_{0}
  \le Ch^{2}\bigl(\Vert u_{t}\Vert_{2}+\Vert p_{t}\Vert_{1}\bigr).
  \label{eq:Stokes-t-est}
\end{equation}

The discrete operator $\Delta_{h}:X^{0}_{h}\to X^{0}_{h}$ is given as follows
\begin{equation}
  (\Delta_{h}\eta_{h},\chi_{h})
  = (\nabla\eta_{h},\nabla\chi_{h}),
  \quad \forall\,\chi_{h}\in X^{0}_{h}.
  \label{eq:Deltah}
\end{equation}
The subspace $V_{h}^{\mathrm{div}}$ of $V_{h}$ is introduced as
\begin{equation}
  V_{h}^{\mathrm{div}}
  = \bigl\{v_{h}\in V_{h} :
    (\operatorname{div}v_{h},q_{h})=0,\;\forall\,q_{h}\in Q_{h}\bigr\}.
  \label{eq:Vhdiv}
\end{equation}
And the discrete Stokes operator $A_{h}:V_{h}^{\mathrm{div}}\to V_{h}^{\mathrm{div}}$
is defined by
\begin{equation}
  (A_{h}u_{h},v_{h}) = (\nabla u_{h},\nabla v_{h}),
  \quad \forall\,v_{h}\in V_{h}^{\mathrm{div}}.
  \label{eq:Ah}
\end{equation}

Then, we have the following Lemmas~\ref{lem:1}--\ref{lem:4}, which can be
found in \cite{Liu2017ch,Diegel2017,Heywood1982,Heywood1990}, or proven
similarly.

\begin{lemma}
  \label{lem:1}
  For $\chi_{h}\in X^{0}_{h}$, there holds
  \begin{equation}
    \Vert\chi_{h}\Vert_{0,\infty}
    + \Vert\nabla\chi_{h}\Vert_{0,6}
    \le C\,\Vert\Delta_{h}\chi_{h}\Vert_{0}.
    \label{eq:lem1}
  \end{equation}
\end{lemma}

\begin{lemma}
  \label{lem:2}
  For $u_{h},v_{h},w_{h}\in V_{h}^{\mathrm{div}}$, there hold
  \begin{align}
    b_{1}(u_{h};v_{h},w_{h})
    &\le C\,\Vert\nabla u_{h}\Vert_{0}
            \Vert\nabla v_{h}\Vert_{0}
            \Vert\nabla w_{h}\Vert_{0},
    \label{eq:b1-est1}\\
    b_{1}(u_{h};v_{h},w_{h})
    &\le C\,\Vert u_{h}\Vert_{0}
            \Vert\nabla v_{h}\Vert_{0}^{1/2}
            \Vert A_{h}v_{h}\Vert_{0}^{1/2}
            \Vert\nabla w_{h}\Vert_{0},
    \label{eq:b1-est2}\\
    b_{1}(u_{h};v_{h},w_{h})
    &\le C\,\Vert\nabla u_{h}\Vert_{0}
            \Vert\nabla v_{h}\Vert_{0}^{1/2}
            \Vert A_{h}v_{h}\Vert_{0}^{1/2}
            \Vert w_{h}\Vert_{0},
    \label{eq:b1-est3}\\
    b_{1}(u_{h};v_{h},w_{h})
    &\le C\,\Vert\nabla u_{h}\Vert_{0}^{1/2}
            \Vert A_{h}u_{h}\Vert_{0}^{1/2}
            \Vert\nabla v_{h}\Vert_{0}
            \Vert w_{h}\Vert_{0},
    \label{eq:b1-est4}\\
    b_{1}(u_{h};v_{h},w_{h})
    &\le C\,\Vert u_{h}\Vert_{0}
            \Vert\nabla u_{h}\Vert_{0}^{1/2}
            \Vert A_{h}v_{h}\Vert_{0}
            \Vert w_{h}\Vert_{0}.
    \label{eq:b1-est5}
  \end{align}
\end{lemma}

\begin{lemma}
  \label{lem:3}
  For $u_{h}\in V_{h}^{\mathrm{div}}$, $\varphi_{h},\chi_{h}\in X^{0}_{h}$,
  there hold
  \begin{align}
    b_{2}(u_{h};\varphi_{h},\chi_{h})
    &\le C\,\Vert\nabla u_{h}\Vert_{0}
            \Vert\nabla\chi_{h}\Vert_{0}
            \Vert\nabla\varphi_{h}\Vert_{0},
    \label{eq:b2-est1}\\
    b_{2}(u_{h};\varphi_{h},\chi_{h})
    &\le C\,\Vert u_{h}\Vert_{0}
            \Vert\nabla\varphi_{h}\Vert_{0}^{1/2}
            \Vert\Delta_{h}\varphi_{h}\Vert_{0}^{1/2}
            \Vert\nabla\chi_{h}\Vert_{0},
    \label{eq:b2-est2}\\
    b_{2}(u_{h};\varphi_{h},\chi_{h})
    &\le C\,\Vert\nabla u_{h}\Vert_{0}^{1/2}
            \Vert A_{h}u_{h}\Vert_{0}^{1/2}
            \Vert\nabla\varphi_{h}\Vert_{0}
            \Vert\chi_{h}\Vert_{0},
    \label{eq:b2-est3}\\
    b_{2}(u_{h};\varphi_{h},\chi_{h})
    &\le C\,\Vert\nabla u_{h}\Vert_{0}
            \Vert\nabla\varphi_{h}\Vert_{0}^{1/2}
            \Vert\Delta_{h}\varphi_{h}\Vert_{0}^{1/2}
            \Vert\chi_{h}\Vert_{0}.
    \label{eq:b2-est4}
  \end{align}
\end{lemma}

\begin{lemma}
  \label{lem:4}
  For $\eta_{h},\varphi_{h},\chi_{h}\in X^{0}_{h}$, there holds
  \begin{equation}
    (\eta_{h}\nabla\varphi_{h},\nabla\chi_{h})
    \le C\,\Vert\nabla\eta_{h}\Vert_{0}^{1/2}
           \Vert\Delta_{h}\eta_{h}\Vert_{0}^{1/2}
           \Vert\Delta_{h}\varphi_{h}\Vert_{0}
           \Vert\chi_{h}\Vert_{0}.
    \label{eq:lem4}
  \end{equation}
\end{lemma}

Finally, assume that the solutions and initial values of equations
\eqref{eq:NSPNP} satisfy the following regularity condition:
\begin{equation}
  \begin{aligned}
    &\Vert u_{0}\Vert_{2}^{2}
     + \Vert c_{10}\Vert_{2}^{2}
     + \Vert c_{20}\Vert_{2}^{2}
     + \Vert c_{1}\Vert_{L^{\infty}(H^{2})}^{2}
     + \Vert c_{2}\Vert_{L^{\infty}(H^{2})}^{2}
     + \Vert\psi\Vert_{L^{\infty}(H^{3})}^{2}\\
    &\quad
     + \Vert u\Vert_{L^{\infty}(H^{2})}^{2}
     + \Vert p\Vert_{L^{\infty}(H^{1})}^{2}
     + \Vert c_{1t}\Vert_{L^{2}(H^{2})}^{2}
     + \Vert c_{1tt}\Vert_{L^{2}(L^{2})}^{2}
     + \Vert c_{2t}\Vert_{L^{2}(H^{2})}^{2}\\
    &\quad
     + \Vert c_{2tt}\Vert_{L^{2}(L^{2})}^{2}
     + \Vert u_{t}\Vert_{L^{2}(H^{2})}^{2}
     + \Vert u_{tt}\Vert_{L^{2}(L^{2})}^{2}
     + \Vert p_{t}\Vert_{L^{2}(H^{1})}^{2}
    \le C_{0}.
  \end{aligned}
  \label{eq:regularity}
\end{equation}

% ------------------------------------------------------------------
\section{The linearized backward Euler scheme}
\label{sec:scheme}
% ------------------------------------------------------------------

Let $\{t_{j}: t_{j}=j\tau,\;0\le j\le J\}$ be a uniform partition of the
time interval $[0,T]$ with time step $\tau=T/J$. Denote
$f^{j}=f(\cdot,t_{j})$, and for a sequence of functions $\{f^{j}\}_{j=0}^{J}$
we define
\begin{equation}
  D_{\tau}f^{j+1} = \frac{f^{j+1}-f^{j}}{\tau}.
  \label{eq:Dtau}
\end{equation}

With the above notation, the linearized backward Euler scheme of
\eqref{eq:variational} is given as follows.

\medskip
\noindent\textbf{Step~1.}
Set $(\tilde{c}^{0}_{1h},\tilde{c}^{0}_{2h},u^{0}_{h})
=(R_{h}\tilde{c}_{10},R_{h}\tilde{c}_{20},S_{h}u_{0})$.
For $j=0,1,\ldots,J-1$, find
$(u^{j+1}_{h},p^{j+1}_{h},\tilde{c}^{j+1}_{1h},
 \tilde{c}^{j+1}_{2h},\psi^{j+1}_{h})
\in V_{h}\times Q_{h}\times X^{0}_{h}\times X^{0}_{h}\times X^{0}_{h}$
such that
\begin{equation}
  \left\{
  \begin{aligned}
    &\bigl(D_{\tau}u^{j+1}_{h},v_{h}\bigr)
     + \bigl(\nabla u^{j+1}_{h},\nabla v_{h}\bigr)
     + b_{1}\bigl(u^{j}_{h};\,u^{j+1}_{h},\,v_{h}\bigr)
     - \bigl(\operatorname{div}v_{h},\,p^{j+1}_{h}\bigr)\\
    &\quad
     = -\bigl((\tilde{c}^{j}_{1h}-\tilde{c}^{j}_{2h})\nabla\psi^{j+1}_{h},\,v_{h}\bigr),
     \quad \forall\,v_{h}\in V_{h},\\[4pt]
    &\bigl(\operatorname{div}u^{j+1}_{h},\,q_{h}\bigr) = 0,
     \quad \forall\,q_{h}\in Q_{h},\\[4pt]
    &\bigl(D_{\tau}\tilde{c}^{j+1}_{1h},\,\eta_{1h}\bigr)
     + \bigl(\nabla\tilde{c}^{j+1}_{1h},\nabla\eta_{1h}\bigr)
     + \bigl((\tilde{c}^{j}_{1h}+m_{0})\nabla\psi^{j+1}_{h},\nabla\eta_{1h}\bigr)\\
    &\quad
     + b_{2}\bigl(u^{j}_{h};\,\tilde{c}^{j+1}_{1h},\,\eta_{1h}\bigr)
     = 0,
     \quad \forall\,\eta_{1h}\in X_{h},\\[4pt]
    &\bigl(D_{\tau}\tilde{c}^{j+1}_{2h},\,\eta_{2h}\bigr)
     + \bigl(\nabla\tilde{c}^{j+1}_{2h},\nabla\eta_{2h}\bigr)
     - \bigl((\tilde{c}^{j}_{2h}+m_{0})\nabla\psi^{j+1}_{h},\nabla\eta_{2h}\bigr)\\
    &\quad
     + b_{2}\bigl(u^{j}_{h};\,\tilde{c}^{j+1}_{2h},\,\eta_{2h}\bigr)
     = 0,
     \quad \forall\,\eta_{2h}\in X_{h},\\[4pt]
    &\bigl(\nabla\psi^{j+1}_{h},\nabla\varphi_{h}\bigr)
     = \bigl(\tilde{c}^{j+1}_{1h}-\tilde{c}^{j+1}_{2h},\,\varphi_{h}\bigr),
     \quad \forall\,\varphi_{h}\in X_{h}.
  \end{aligned}
  \right.
  \label{eq:scheme}
\end{equation}

\noindent\textbf{Step~2.} Set
\begin{equation}
  c^{j+1}_{1h} = \tilde{c}^{j+1}_{1h} + m_{0},
  \qquad
  c^{j+1}_{2h} = \tilde{c}^{j+1}_{2h} + m_{0}.
  \label{eq:c-recover}
\end{equation}

To derive the well-posedness and unconditional error estimates of scheme
\eqref{eq:scheme}, we need the following inductive hypothesis:
\begin{equation}
  \bigl\Vert\tilde{c}^{j}_{1h}\bigr\Vert_{0,3}
  + \bigl\Vert\tilde{c}^{j}_{2h}\bigr\Vert_{0,3}
  \le C_{1}.
  \label{eq:inductive}
\end{equation}

In the remainder of this section we prove the existence and uniqueness of
scheme \eqref{eq:scheme}.

\begin{theorem}
  \label{thm:wellposed}
  If the time step $\tau$ is small enough, scheme \eqref{eq:scheme} has a
  unique solution.
\end{theorem}

\begin{proof}
Since scheme \eqref{eq:scheme} is an algebraic linear system, we only need
to prove uniqueness. Suppose there exist
$(\tilde{c}^{j+1}_{1h,i},\tilde{c}^{j+1}_{2h,i},\psi^{j+1}_{h,i})
\in X^{0}_{h}\times X^{0}_{h}\times X^{0}_{h}$, $i=1,2$, satisfying the
last three equations of \eqref{eq:scheme}. Then we have
\begin{equation}
  \left\{
  \begin{aligned}
    &\bigl(D_{\tau}(\tilde{c}^{j+1}_{1h,1}-\tilde{c}^{j+1}_{1h,2}),\,\eta_{1h}\bigr)
     + \bigl(\nabla(\tilde{c}^{j+1}_{1h,1}-\tilde{c}^{j+1}_{1h,2}),\nabla\eta_{1h}\bigr)\\
    &\quad
     + \bigl((\tilde{c}^{j}_{1h}+m_{0})
       \nabla(\psi^{j+1}_{h,1}-\psi^{j+1}_{h,2}),\nabla\eta_{1h}\bigr)
     + b_{2}\bigl(u^{j}_{h};\,
       \tilde{c}^{j+1}_{1h,1}-\tilde{c}^{j+1}_{1h,2},\,\eta_{1h}\bigr) = 0,\\[4pt]
    &\bigl(D_{\tau}(\tilde{c}^{j+1}_{2h,1}-\tilde{c}^{j+1}_{2h,2}),\,\eta_{2h}\bigr)
     + \bigl(\nabla(\tilde{c}^{j+1}_{2h,1}-\tilde{c}^{j+1}_{2h,2}),\nabla\eta_{2h}\bigr)\\
    &\quad
     - \bigl((\tilde{c}^{j}_{2h}+m_{0})
       \nabla(\psi^{j+1}_{h,1}-\psi^{j+1}_{h,2}),\nabla\eta_{2h}\bigr)
     + b_{2}\bigl(u^{j}_{h};\,
       \tilde{c}^{j+1}_{2h,1}-\tilde{c}^{j+1}_{2h,2},\,\eta_{2h}\bigr) = 0,\\[4pt]
    &\bigl(\nabla(\psi^{j+1}_{h,1}-\psi^{j+1}_{h,2}),\nabla\varphi_{h}\bigr)
     = \bigl((\tilde{c}^{j+1}_{1h,1}-\tilde{c}^{j+1}_{1h,2})
       - (\tilde{c}^{j+1}_{2h,1}-\tilde{c}^{j+1}_{2h,2}),\,\varphi_{h}\bigr).
  \end{aligned}
  \right.
  \label{eq:diff-system}
\end{equation}

Taking $\varphi_{h}=\psi^{j+1}_{h,1}-\psi^{j+1}_{h,2}$ and
$\varphi_{h}=\Delta_{h}(\psi^{j+1}_{h,1}-\psi^{j+1}_{h,2})$ in the third
equation of \eqref{eq:diff-system}, respectively, we obtain
\begin{equation}
  \bigl\Vert\nabla(\psi^{j+1}_{h,1}-\psi^{j+1}_{h,2})\bigr\Vert_{0}
  \le C\bigl(\bigl\Vert\tilde{c}^{j+1}_{1h,1}-\tilde{c}^{j+1}_{1h,2}\bigr\Vert_{0}
           + \bigl\Vert\tilde{c}^{j+1}_{2h,1}-\tilde{c}^{j+1}_{2h,2}\bigr\Vert_{0}\bigr),
  \label{eq:psi-diff1}
\end{equation}
\begin{equation}
  \bigl\Vert\Delta_{h}(\psi^{j+1}_{h,1}-\psi^{j+1}_{h,2})\bigr\Vert_{0}
  \le \bigl\Vert\tilde{c}^{j+1}_{1h,1}-\tilde{c}^{j+1}_{1h,2}\bigr\Vert_{0}
    + \bigl\Vert\tilde{c}^{j+1}_{2h,1}-\tilde{c}^{j+1}_{2h,2}\bigr\Vert_{0}.
  \label{eq:psi-diff2}
\end{equation}

Setting $\eta_{1h}=\tilde{c}^{j+1}_{1h,1}-\tilde{c}^{j+1}_{1h,2}$ in the
first equation of \eqref{eq:diff-system} and using
\eqref{eq:lem1}, \eqref{eq:inductive}, \eqref{eq:psi-diff2}, we arrive at
\begin{equation}
  \bigl\Vert\tilde{c}^{j+1}_{1h,1}-\tilde{c}^{j+1}_{1h,2}\bigr\Vert_{0}^{2}
  + \tfrac{\tau}{2}\bigl\Vert\nabla(\tilde{c}^{j+1}_{1h,1}-\tilde{c}^{j+1}_{1h,2})\bigr\Vert_{0}^{2}
  \le CC_{1}^{2}\tau\bigl(\bigl\Vert\tilde{c}^{j+1}_{1h,1}-\tilde{c}^{j+1}_{1h,2}\bigr\Vert_{0}^{2}
      + \bigl\Vert\tilde{c}^{j+1}_{2h,1}-\tilde{c}^{j+1}_{2h,2}\bigr\Vert_{0}^{2}\bigr).
  \label{eq:c1-unique}
\end{equation}

Similarly, we get
\begin{equation}
  \bigl\Vert\tilde{c}^{j+1}_{2h,1}-\tilde{c}^{j+1}_{2h,2}\bigr\Vert_{0}^{2}
  + \tfrac{\tau}{2}\bigl\Vert\nabla(\tilde{c}^{j+1}_{2h,1}-\tilde{c}^{j+1}_{2h,2})\bigr\Vert_{0}^{2}
  \le CC_{1}^{2}\tau\bigl(\bigl\Vert\tilde{c}^{j+1}_{1h,1}-\tilde{c}^{j+1}_{1h,2}\bigr\Vert_{0}^{2}
      + \bigl\Vert\tilde{c}^{j+1}_{2h,1}-\tilde{c}^{j+1}_{2h,2}\bigr\Vert_{0}^{2}\bigr).
  \label{eq:c2-unique}
\end{equation}

Summing \eqref{eq:c1-unique} and \eqref{eq:c2-unique} gives
\begin{equation}
  \begin{aligned}
    &(1-CC_{1}^{2}\tau)
     \bigl(\bigl\Vert\tilde{c}^{j+1}_{1h,1}-\tilde{c}^{j+1}_{1h,2}\bigr\Vert_{0}^{2}
           + \bigl\Vert\tilde{c}^{j+1}_{2h,1}-\tilde{c}^{j+1}_{2h,2}\bigr\Vert_{0}^{2}\bigr)\\
    &\quad
     + \tfrac{\tau}{2}\bigl(\bigl\Vert\nabla(\tilde{c}^{j+1}_{1h,1}-\tilde{c}^{j+1}_{1h,2})\bigr\Vert_{0}^{2}
           + \bigl\Vert\nabla(\tilde{c}^{j+1}_{2h,1}-\tilde{c}^{j+1}_{2h,2})\bigr\Vert_{0}^{2}\bigr)
     \le 0.
  \end{aligned}
  \label{eq:c12-sum}
\end{equation}

Setting $\tau<\tau_{1}:=\dfrac{1}{2CC_{1}^{2}}$ in \eqref{eq:c12-sum} yields
\begin{equation}
  \tilde{c}^{j+1}_{1h,1} = \tilde{c}^{j+1}_{1h,2},
  \qquad
  \tilde{c}^{j+1}_{2h,1} = \tilde{c}^{j+1}_{2h,2},
  \label{eq:c-unique}
\end{equation}
which together with \eqref{eq:psi-diff1} leads to
\begin{equation}
  \psi^{j+1}_{h,1} = \psi^{j+1}_{h,2}.
  \label{eq:psi-unique}
\end{equation}
Finally, the first two equations of \eqref{eq:scheme} represent the
linearized backward Euler scheme of the Navier--Stokes equations, which has
the unique solution $(u^{j+1}_{h},p^{j+1}_{h})$. The proof is completed.
\end{proof}

% ============================================================
%  BIBLIOGRAPHY ENTRIES (representative subset)
%  Include in your .bib file or use \bibliography{refs}
% ============================================================
%
% @book{John2016,
%   author    = {John, Volker},
%   title     = {Finite Element Methods for Incompressible Flow Problems},
%   publisher = {Springer-Verlag},
%   address   = {Berlin},
%   year      = {2016}
% }
%
% @article{Li2021,
%   author  = {Li, Zhenzhen and Li, Minghao and Shi, Dongyang},
%   title   = {Unconditional convergence and superconvergence analysis for the
%              transient {S}tokes equations with damping},
%   journal = {Appl. Math. Comput.},
%   volume  = {389},
%   pages   = {125572},
%   year    = {2021}
% }
%
% @article{Liu2017ch,
%   author  = {Liu, Yali and Chen, Wenbin and Wang, Cheng and Wise, Steven M.},
%   title   = {Error analysis of a mixed finite element method for a
%              {C}ahn--{H}illiard--{H}ele--{S}haw system},
%   journal = {Numer. Math.},
%   volume  = {135},
%   pages   = {679--709},
%   year    = {2017}
% }
%
% @article{Diegel2017,
%   author  = {Diegel, Amanda E. and Wang, Cheng and Wang, Xiaoming and Wise, Steven M.},
%   title   = {Convergence analysis and error estimates for a second order accurate
%              finite element method for the {C}ahn--{H}illiard--{N}avier--{S}tokes system},
%   journal = {Numer. Math.},
%   volume  = {137},
%   pages   = {495--534},
%   year    = {2017}
% }
%
% @article{Heywood1982,
%   author  = {Heywood, John G. and Rannacher, Rolf},
%   title   = {Finite element approximation of the nonstationary {N}avier--{S}tokes
%              problem. {I}. {R}egularity of solutions and second-order error estimates
%              for spatial discretization},
%   journal = {SIAM J. Numer. Anal.},
%   volume  = {19},
%   number  = {2},
%   pages   = {275--311},
%   year    = {1982}
% }
%
% @article{Heywood1990,
%   author  = {Heywood, John G. and Rannacher, Rolf},
%   title   = {Finite-element approximation of the nonstationary {N}avier--{S}tokes
%              problem. {P}art {IV}: {E}rror analysis for second-order time discretization},
%   journal = {SIAM J. Numer. Anal.},
%   volume  = {27},
%   number  = {2},
%   pages   = {353--384},
%   year    = {1990}
% }
